% -*- coding: utf-8 -*-
%%%%%%%%%%%%%%%%%%%%%%%%%%%%%%%%%%%%%%%%%%%%%%%%%%%%%%%%%%%
% 主文档 格式定义
%%%%%%%%%%%%%%%%%%%%%%%%%%%%%%%%%%%%%%%%%%%%%%%%%%%%%%%%%%%

% 将版芯控制在240mm以内, 正文范围控制在220mm以内
% 将版芯控制在240mm以内, 正文范围控制在220mm以内
\addtolength{\headsep}{-10pt}          %页眉位置
%\addtolength{\footskip}{0}         %页脚位置

%%%%%%%%%%%%%%%%%%%%%%%%%%%%%%%%%%%%%%%%%%%%%%%%%%%%%%%%%%
%  页面设置
%%%%%%%%%%%%%%%%%%%%%%%%%%%%%%%%%%%%%%%%%%%%%%%%%%%%%%%%%
%\setlength{\voffset}{-12pt}
%\setlength{\topmargin}{12pt}
%\addtolength{\textheight}{22pt}
%
%\setlength{\hoffset}{15pt}
%\setlength{\oddsidemargin}{20pt}
%\setlength{\evensidemargin}{20pt}

%%%%%%%%%%%%%%%%%%%%%%%%%%%%%%%%%%%%%%%%%%%%%%%%%%%%%%%%%%%
% 公式的精调
%%%%%%%%%%%%%%%%%%%%%%%%%%%%%%%%%%%%%%%%%%%%%%%%%%%%%%%%%%%
%\setlength{\mathindent}{4.7 em}     %左对齐公式缩进量

% \eqnarray如果很长，影响分栏、换行和分页（整块挪动，造成页面空白），
% 可以设置成为自动调整模式
\allowdisplaybreaks[4]

%%%%%%%%%%%%%%%%%%%%%%%%%%%%%%%%%%%%%%%%%%%%%%%%%%%%%%%%%%%
%下面这组命令使浮动对象的缺省值稍微宽松一点，从而防止幅度
%对象占据过多的文本页面，也可以防止在很大空白的浮动页上放置
%很小的图形。
%%%%%%%%%%%%%%%%%%%%%%%%%%%%%%%%%%%%%%%%%%%%%%%%%%%%%%%%%%%
\renewcommand{\textfraction}{0.15}
\renewcommand{\topfraction}{0.85}
\renewcommand{\bottomfraction}{0.65}
\renewcommand{\floatpagefraction}{0.60}


%%%%%%%%%%%%%%%%%%%%%%%%%%%%%%%%%%%%%%%%%%%%%%%%%%%%%%%%%%%
%下面这组命令可以使公式编号随着每开始新的一节而重新开始。
%%%%%%%%%%%%%%%%%%%%%%%%%%%%%%%%%%%%%%%%%%%%%%%%%%%%%%%%%%%
%\makeatletter      % '@' is now a normail "letter" for TeX
%\@addtoreset{eqation}{section}
%\makeatother       % '@' is restored as a "non-letter" character for TeX

%%%%%%%%%%%%%%%%%%%%%%%%%%%%%%%%%%%%%%%%%%%%%%%%%%%%%%%%%%%
% 重定义字体命令
%%%%%%%%%%%%%%%%%%%%%%%%%%%%%%%%%%%%%%%%%%%%%%%%%%%%%%%%%%%
% 注意win2000没有 simsun,　最好到网上找一个
% 一些字体是office2000 带的
%%%%%%%%%%%%%%%%%%%%%%%%%%%%%%%%%%%%%%%%%%%%%%%%%%%%%%%%%%%
%
\newcommand{\song}{\CJKfamily{song}}    % 宋体   (Windows自带simsun.ttf)
\newcommand{\fs}{\CJKfamily{fs}}        % 仿宋体 (Windows自带simfs.ttf)
\newcommand{\kai}{\CJKfamily{kai}}      % 楷体   (Windows自带simkai.ttf)
\newcommand{\hei}{\CJKfamily{hei}}      % 黑体   (Windows自带simhei.ttf)
\newcommand{\li}{\CJKfamily{li}}        % 隶书   (Windows自带simli.ttf)

%%%%%%%%%%%%%%%%%%%%%%%%%%%%%%%%%%%%%%%%%%%%%%%%%%%%%%%%%%%
% 重定义字号命令
%%%%%%%%%%%%%%%%%%%%%%%%%%%%%%%%%%%%%%%%%%%%%%%%%%%%%%%%%%%
%
\newcommand{\yihao}{\zihao{1}}    % 一号, 1.4倍行距
\newcommand{\erhao}{\zihao{2}}    % 二号, 1.25倍行距
\newcommand{\xiaoer}{\zihao{-2}}    % 小二, 单倍行距
\newcommand{\sanhao}{\zihao{3}}    % 三号, 1.5倍行距
\newcommand{\xiaosan}{\zihao{-3}}    % 小三, 1.5倍行距
\newcommand{\sihao}{\zihao{4}}    % 四号, 1.5倍行距
\newcommand{\banxiaosi}{\fontsize{13pt}{19.5pt}\selectfont}    % 半小四, 1.5倍行距
\newcommand{\xiaosi}{\zihao{-4}}    % 小四, 1.5倍行距
\newcommand{\dawuhao}{\fontsize{11pt}{11pt}\selectfont}    % 大五号, 单倍行距
\newcommand{\wuhao}{\zihao{5}}    % 五号, 单倍行距


%%%%%%%%%%%%%%%%%%%%%%%%%%%%%%%%%%%%%%%%%%%%%%%%%%%%%%%%%%%
% 重定义一些正文相关标题
%%%%%%%%%%%%%%%%%%%%%%%%%%%%%%%%%%%%%%%%%%%%%%%%%%%%%%%%%%%

%\theoremstyle{plain} \theorembodyfont{\song\rmfamily}
%\theoremheaderfont{\hei\rmfamily} \theoremseparator{:}
%
%\renewcommand\contentsname{目~~~~~~录}
%\renewcommand\listfigurename{插图}
%\renewcommand\listtablename{表格}
%\renewcommand\indexname{索引}
%\renewcommand\figurename{图}
%\renewcommand\tablename{表}

%\newtheorem{definition}{\hei 定义}[chapter]
%\newtheorem{proposition}[definition]{\hei 命题}
%\newtheorem{lemma}[definition]{\hei 引理}
%\newtheorem{theorem}{\hei 定理}[chapter]
%\newtheorem{axiom}{\hei 公理}
%\newtheorem{corollary}[definition]{\hei 推论}
%\newtheorem{exercise}[definition]{}

%\theoremheaderfont{\CJKfamily{hei}\rmfamily}\theorembodyfont{\rmfamily}
%\theoremstyle{nonumberplain} \theoremseparator{:}
%\theoremsymbol{$\blacksquare$}
%\newtheorem{proof}{\hei 证明}

%\theoremsymbol{$\square$}
%\newtheorem{example}{\hei 例}


%%%%%%%%%%%%%%%%%%%%%%%%%%%%%%%%%%%%%%%%%%%%%%%%%%%%%%%
% 定义页眉和页脚 使用fancyhdr 宏包
%%%%%%%%%%%%%%%%%%%%%%%%%%%%%%%%%%%%%%%%%%%%%%%%%%%%%%%%
%\newcommand{\makeheadrule}{%
%    \makebox[0pt][l]{\rule[.7\baselineskip]{\headwidth}{0.8pt}}%
%    \vskip-.8\baselineskip}%
%
%  \makeatletter
%  \renewcommand{\headrule}{%
%    {\if@fancyplain\let\headrulewidth\plainheadrulewidth\fi
%      \makeheadrule}}%
%
%  \pagestyle{fancy}
%
%去掉章节标题中的数字
%\renewcommand{\chaptermark}[1]{\markboth{\chaptername \ #1}{}}

%\fancyhf{}
%\fancyhead[CO]{\wuhao\leftmark}
%\fancyhead[CE]{\wuhao\leftmark}
%\fancyfoot[C,C]{\wuhao\thepage}

%%%%%%%%%%%%%%%%%%%%%%%%%%%%%%%%%%%%%%%%%%%%%%%%%%%%%%%%
% 设置行距和段落间垂直距离
%%%%%%%%%%%%%%%%%%%%%%%%%%%%%%%%%%%%%%%%%%%%%%%%%%%%%%%%

% 段落之间的竖直距离
\setlength{\parskip}{3pt plus1pt minus1pt}

% 定义行距
\renewcommand{\baselinestretch}{1.2}


%%%%%%%%%%%%%%%%%%%%%%%%%%%%%%%%%%%%%%%%%%%%%%%%%%%%%%%
% 修改引用的格式，
%%%%%%%%%%%%%%%%%%%%%%%%%%%%%%%%%%%%%%%%%%%%%%%%%%%%%%%
% 增加 \ucite 命令使显示的引用为上标形式
\newcommand{\ucite}[1]{$^{\mbox{\scriptsize \cite{#1}}}$}

%%%%%%%%%%%%%%%%%%%%%%%%%%%%%%%%%%%%%%%%%%%%%%%%%%%%%%%%%%%
% 定制浮动图形和表格标题样式
%%%%%%%%%%%%%%%%%%%%%%%%%%%%%%%%%%%%%%%%%%%%%%%%%%%%%%%%%%%
%添加新的表格控制符，用法与p{xx}相似
%P{}表示定制表格例的宽度并且使不够长度的居中与同等其他列上对齐
%M{}表示定制表格例的宽度并且使不够长度的居中与同等其他列中对齐
%B{}表示定制表格例的宽度并且使不够长度的居中与同等其他列下对齐
%\newcommand{\PreserveBackslash}[1]{\let\temp=\\#1\let\\=\temp}
%\newcolumntype{P}[1]{>{\PreserveBackslash\centering}p{#1}}
%\newcolumntype{M}[1]{>{\PreserveBackslash\centering}m{#1}}
%\newcolumntype{B}[1]{>{\PreserveBackslash\centering}b{#1}}

\renewcommand{\captionfont}{\CJKfamily{song}\rmfamily}
\renewcommand{\captionlabelfont}{\CJKfamily{song}\rmfamily}

% 按清华标准, 去掉图表号后面的:
%不太清楚南大的标准，所以没改
\captionsetup{labelsep=blank}

% 按清华标准, 图表标题字体为11pt, 这里写作大五号
\renewcommand{\captionfont}{\dawuhao}

% 设置Caption的宽度
\setcaptionmargin{.5in}

% define graph scale
\def\gs{1.0}

%%%%%%%%%%%%%%%%%%%%%%%%%%%%%%%%%%%%%%%%%%%%%%%%%%%%%%%%%%%%%%%%%%%%%%
% 自定义项目列表标签及格式 \begin{denselist} 列表项 \end{denselist}
%%%%%%%%%%%%%%%%%%%%%%%%%%%%%%%%%%%%%%%%%%%%%%%%%%%%%%%%%%%%%%%%%%%%%%
\newcounter{newlist} %自定义新计数器
\newenvironment{denselist}[1][可改变的列表题目]{%%%%%定义新环境
\begin{list}{\textbf{\hei #1} \arabic{newlist}：} %%标签格式
    {
    \usecounter{newlist}
     \setlength{\labelwidth}{22pt} %标签盒子宽度
     \setlength{\labelsep}{0cm} %标签与列表文本距离
     \setlength{\leftmargin}{0cm} %左右边界
     \setlength{\rightmargin}{0cm}
     \setlength{\parsep}{0ex} %段落间距
     \setlength{\itemsep}{0ex} %标签间距
     \setlength{\itemindent}{44pt} %标签缩进量
     \setlength{\listparindent}{22pt} %段落缩进量
    }}
{\end{list}}%%%%%

\def\defaultfont{\renewcommand{\baselinestretch}{1.5}
\fontsize{12pt}{13pt}\selectfont}

%%%%%%%%%%%%%%%%%%%%%%%%%%%%%%%%%%%%%%%%%%%%%%%%%%%%%%%%%%%%%%%%%%%%%%
% 封面、摘要、与目录格式定义 --- by Guanghong Zuo Feng Hua & Lei Wang
%%%%%%%%%%%%%%%%%%%%%%%%%%%%%%%%%%%%%%%%%%%%%%%%%%%%%%%%%%%%%%%%%%%%%%
\makeatletter
\def\ctitle#1{\def\@ctitle{#1}}\def\@ctitle{}
\def\cdegree#1{\def\@cdegree{#1}}\def\@cdegree{}
\def\cdepart#1{\def\@cdepart{#1}}\def\@cdepart{}
\def\csubject#1{\def\@csubject{#1}}\def\@csubject{}
\def\cdirection#1{\def\@cdirection{#1}}\def\@cdirection{}
\def\cauthor#1{\def\@cauthor{#1}}\def\@cauthor{}
\def\csupervisor#1{\def\@csupervisor{#1}}\def\@csupervisor{}
\def\stitle#1{\def\@stitle{#1}}\def\@stitle{}
\def\cdate#1{\def\@cdate{#1}}\def\@cdate{}
\long\def\cabstract#1{\long\def\@cabstract{#1}}\long\def\@cabstract{}
\def\ckeywords#1{\def\@ckeywords{#1}}\def\@ckeywords{}

\def\etitle#1{\def\@etitle{#1}}\def\@etitle{}
\def\edegree#1{\def\@edegree{#1}}\def\@edegree{}
\def\edepart#1{\def\@edepart{#1}}\def\@edepart{}
\def\esubject#1{\def\@esubject{#1}}\def\@esubject{}
\def\edirection#1{\def\@edirection{#1}}\def\@edirection{}
\def\eauthor#1{\def\@eauthor{#1}}\def\@eauthor{}
\def\esupervisor#1{\def\@esupervisor{#1}}\def\@esupervisor{}
\def\estitle#1{\def\@estitle{#1}}\def\@estitle{}
\def\eaddress#1{\def\@eaddress{#1}}\def\@eaddress{}
\def\edate#1{\def\@edate{#1}}\def\@edate{}
\long\def\eabstract#1{\long\def\@eabstract{#1}}\long\def\@eabstract{}
\def\ekeywords#1{\def\@ekeywords{#1}}\def\@ekeywords{}

\def\ngrade#1{\def\@ngrade{#1}}\def\@ngrade{}
\def\snumber#1{\def\@snumber{#1}}\def\@snumber{}

\sloppy
\pagenumbering{roman}

\def\makecover{
  \begin{titlepage}
\includegraphics[width=2.5cm]{njuxh}\\
    \parbox[t][.5cm][t]{\textwidth}{
    \begin{center}  \end{center} }
  \begin{center}

      \parbox[t][5.2cm][t]{\textwidth}{\begin{center} \song
          \includegraphics[width=5.0cm]{njuwz}\\
         {\bf\yihao 本~~~科~~~毕~~~业~~~论~~~文\\}
         \end{center}}

    \parbox[t][3.5cm][t]{\textwidth}{
    \begin{center}  \end{center} }


    \parbox[t][6cm][c]{\textwidth}{\begin{center}
      \fontsize{16pt}{8pt}\selectfont \song
      \begin{tabular}{p{4.5em} p{2em} p{4em}@{} p{2.5em} p{6em}}
    院~~~~~~~~系 & \multicolumn{4}{c}{\@cdepart}\\
    \cmidrule(r){2-5}\addlinespace[10pt]
    专~~~~~~~~业 & \multicolumn{4}{c}{\@csubject}\\
    \cmidrule(r){2-5}\addlinespace[10pt]
    题~~~~~~~~目 &  \multicolumn{4}{c}{\@ctitle}\\
    \cmidrule(r){2-5}\addlinespace[10pt]
    年~~~~~~~~级 &  \multicolumn{2}{c}{\@ngrade} & 学号 & \@snumber\\
    \cmidrule(r){2-3} \cmidrule(r){5-5}\addlinespace[10pt]
    学生姓名 & \multicolumn{4}{c}{\@cauthor}\\
    \cmidrule(r){2-5}\addlinespace[10pt]
    指导教师 & \multicolumn{2}{c}{\@csupervisor} & 职称 & \@stitle\\
    \cmidrule(r){2-3} \cmidrule(r){5-5}\addlinespace[10pt]
    \multicolumn{2}{l}{论文提交日期} & \multicolumn{3}{c}{\@cdate}\\
    \cmidrule(r){3-5}
  \end{tabular}\end{center}}

    \parbox[t][2.5cm][t]{\textwidth}{
    \begin{center}  \end{center} }

    %\parbox[t][1cm][b]{\textwidth}{
    %\begin{center} {\sanhao \kai \@cdate} \end{center} }

  \end{center}

    % 封二
    \newpage

    \thispagestyle{empty}

    \parbox[t][16cm][t]{\textwidth}{
    \begin{center}  \end{center} }

    \parbox[b][10em][b]{\textwidth}{\kai \sanhao
      \begin{tabular}{l p{12em}}
        学~~~~~~~~~~~~~~~~号：& {\bf ~~~~\@snumber}\\
        论文答辩日期：& ~~~~~~~~~~年~~~~~~~月~~~~~~~日\\
        指~~~导~~~教~~~师： &~~~~~~~~~~~~~~~~~~~~~~（签字）\\
        \end{tabular}
    }

    % 英文封面
    \newpage

    \thispagestyle{empty}
    \begin{center}
    \parbox[t][1cm][t]{\textwidth}{
    \begin{center}  \end{center} }

    \parbox[t][1cm][c]{\textwidth}{\sanhao
    \begin{center} {\textbf{\textsf Dissertation for {\@edegree}}}\end{center} }

    \parbox[t][4cm][c]{\textwidth}{\erhao
    \begin{center} {\textbf{\textsf{\@etitle}}}\end{center} }

    \parbox[t][4cm][t]{\textwidth}{
    \begin{center}  \end{center} }

    \parbox[t][5cm][b]{\textwidth}{
    \begin{center} {\sanhao \textsf{\textbf{\@eauthor\\
            Supervised by ~\@esupervisor \\ \@esubject\\} }} \end{center} }

    \parbox[t][3cm][c]{\textwidth}{ {\sanhao
    \begin{center}
      \@eaddress
    \end{center} } }

    \parbox[t][2cm][b]{\textwidth}{
    \begin{center} {\sanhao \textbf{\textsf{\@edate}}} \end{center} }
    \end{center}
    \end{titlepage}

    \clearpage

    %Abstract%%%%%%%%%%%%%%%%%%%%%%%%%%%%%%%%%%%%%%%%%%%%%%%%%%%%%
    \defaultfont
    \pagestyle{plain}

      \phantomsection
      \begin{center} \kai \xiaoer
        \uuline{南京大学本科生毕业论文(设计)中文摘要}

        \vspace{15pt}

        {\fontsize{16pt}{8pt}\selectfont
          \begin{tabular}{p{3.5em}@{} p{2.5em} p{2em}@{} p{2em}@{} p{2.5em} p{2em}@{} p{7em}@{} p{3.5em}}
            \multicolumn{3}{l}{毕业论文题目:} &\multicolumn{5}{b{17em}}{\@ctitle}\\
            \cmidrule(r){4-8}\addlinespace[8pt]
            \@cdepart &院系& \multicolumn{2}{c}{\@csubject} &专业& \@ngrade &级本科生姓名:& \@cauthor\\
            \cmidrule(r){1-1}\cmidrule(r){3-4}\cmidrule(r){6-6}\cmidrule(r){8-8}\addlinespace[8pt]
            \multicolumn{4}{c}{指导教师(姓名、职称):} &\multicolumn{2}{c}{\@csupervisor} &\@stitle\\
            \cmidrule{5-8}
          \end{tabular}}

        \vspace{1.5cm}

      % {\hei \sanhao 中文摘要}
      \end{center}
      \addcontentsline{toc}{chapter}{中文摘要}
      \markboth{}{}
      \setcounter{page}{1}
      {\kai \sanhao 摘要：} \\

      \@cabstract\\
      \\
     {\kai \sanhao 关键词：} ~\@ckeywords
 
      \defaultfont
      \clearpage

      \newpage
      \phantomsection
      \begin{center} \kai \xiaoer
        \uuline{南京大学本科生毕业论文(设计)英文摘要}
      \end{center}

      \vspace{-10pt}

      {\xiaosan \parindent 0pt

        \begin{tabular}{@{}l p{22em}}
          THESIS: & \fontsize{16pt}{10pt}\selectfont \@etitle\\
        \end{tabular}

        DEPARTMENT:~~~~\@edepart
        \vspace{-5pt}

        SPECIALIZATION:~~~~\@esubject
        \vspace{5pt}

        UNDERGRADUATE:~~~~\@eauthor
        \vspace{5pt}

        MENTOR:~~~~\@esupervisor
        \vspace{5pt}

        ABSTRACT:~~~~
      }

      \vspace{1cm}

      %\begin{center}{\bf \sanhao Abstract}\end{center}
      \addcontentsline{toc}{chapter}{英文摘要}
      \markboth{}{}

      \@eabstract\\
      \\
      KEYWORDS:~\@ekeywords

      \clearpage

    % Contents %%%%%%%%%%%%%%%%%%%%%%%%%%%%%%%%%%%%%%%%%%%%%%%%%%%%%%
    %\renewcommand{\baselinestretch}{1.5}
    %\fontsize{12pt}{12pt}\selectfont
    %\pagestyle{fancy}

    % \phantomsection
    % \addcontentsline{toc}{chapter}{\contentsname}
    \tableofcontents
    \markboth{\contentsname}{\contentsname}
    \clearpage
    %\defaultfont
    \pagenumbering{arabic}
    \setcounter{page}{1}

}

\makeatother
